\documentclass{beamer}
\usepackage{goyabeamer}

% ==========================================
% CONFIGURACIÓN ESTÉTICA (Estilo Metropolis)
% ==========================================

\usepackage{amssymb}
\usepackage{tikz}
\usetikzlibrary{shapes.geometric, arrows, positioning, shadows}

% Estilos TikZ para diagramas profesionales (nanobanana-style)
\tikzset{
    node/.style={circle, minimum width=0.8cm, minimum height=0.8cm, text centered, draw=myblue, fill=mygreen!20, font=\tiny, drop shadow},
    arrow/.style={thick,->,>=stealth, myblue},
    cloud/.style={draw, ellipse, fill=mygray, node distance=2.5cm, minimum height=2em}
}

% ==========================================
% PORTADA
% ==========================================
\title{Bases de Datos No Relacionales (RA7)}
\subtitle{Sistemas Distribuidos y Ecosistema NoSQL}
\author{Gestión de Bases de Datos}
\date{\today}
\institute{FP Madrid - DAM / DAW}

\usepackage[spanish, es-noshorthands, es-tabla]{babel}
\usepackage[autostyle]{csquotes}
\begin{document}

\maketitle

% --- INCLUSIÓN DE SECCIONES ---
\section{1. Motivación: ¿Por qué NoSQL?}

\begin{frame}{El contexto: La era del SQL (1970 - 2000)}
    Durante décadas, las bases de datos relacionales (MySQL, Oracle, SQL Server) fueron la única opción.
    \begin{itemize}
        \item \textbf{Foco:} Ahorrar espacio en disco (que era muy caro).
        \item \textbf{Técnica:} Normalización (dividir datos en muchas tablas para no repetir nada).
        \item \textbf{Garantía:} ACID (Consistencia absoluta).
    \end{itemize}
    \begin{exampleblock}{¿Qué cambió?}
        La llegada de la Web 2.0 y el Big Data. El problema ya no es el espacio en disco, sino la \textbf{velocidad} y el \textbf{volumen} de los datos.
    \end{exampleblock}
\end{frame}

\begin{frame}{Caso 1: Redes Sociales (Complejidad)}
  \textbf{El problema:} El coste de los JOINs.
  \begin{itemize}
    \item En SQL, para saber qué le gusta a los amigos de ``Juan'', el motor debe cruzar la tabla \texttt{Usuarios}, \texttt{Amistades} y \texttt{Likes}.
    \item Cada \texttt{JOIN} es una operación matemática costosa.
    \item \textbf{Recursividad:} ``¿Qué recomiendan los amigos de mis amigos?''. En SQL esto requiere consultas anidadas que bloquean el servidor.
  \end{itemize}
  \begin{alertblock}{Rendimiento Exponencial}
    En SQL, cuanta más relación hay entre los datos, más lenta es la consulta. NoSQL (Grafos) soluciona esto con punteros físicos entre nodos.
  \end{alertblock}
\end{frame}

\begin{frame}[fragile]{Caso 2: Esquemas Rígidos y EAV}
  \textbf{El problema:} ¿Qué pasa si cada fila tiene columnas distintas?
  En SQL usamos el modelo \textbf{EAV (Entity-Attribute-Value)} para simular flexibilidad:
  \begin{lstlisting}[language=SQL]
-- Para sacar un producto con 3 atributos:
SELECT p.nombre, a1.valor as color, a2.valor as talla
FROM Productos p
JOIN Atributos a1 ON p.id = a1.p_id AND a1.attr = 'Color'
JOIN Atributos a2 ON p.id = a2.p_id AND a2.attr = 'Talla';
  \end{lstlisting}
  \begin{itemize}
    \item \textbf{Efecto:} El código se vuelve una ``sopa de JOINs'' ilegible.
    \item \textbf{NoSQL:} Simplemente guardamos un JSON con los campos que necesitemos para esa fila concreta.
  \end{itemize}
\end{frame}

\begin{frame}{Caso 3: Escrituras Masivas (IoT/Logs)}
  \textbf{El problema:} Los bloqueos (Locks).
  \begin{itemize}
    \item SQL es ``celoso''. Para asegurar que nadie lea un dato mientras se escribe, bloquea filas o tablas enteras.
    \item Si tienes 10.000 sensores enviando datos por segundo, la base de datos pasa más tiempo gestionando bloqueos que escribiendo.
  \end{itemize}
  \begin{exampleblock}{Escritura secuencial}
    Sistemas como Cassandra no ``editan'' datos, solo ``añaden'' al final de un archivo. Es como escribir en una libreta sin borrar nunca; es mucho más rápido.
  \end{exampleblock}
\end{frame}

\begin{frame}{Caso 4: Escalado Global (Latencia)}
  \textbf{El problema:} La física y la velocidad de la luz.
  \begin{itemize}
    \item Un paquete de red tarda unos 200ms en ir y volver de Madrid a Tokio.
    \item En SQL (ACID), una transacción no termina hasta que todos los nodos confirman.
    \item \textbf{Resultado:} Tu aplicación en Madrid va lenta porque el servidor de Tokio tarda en contestar.
  \end{itemize}
  \begin{alertblock}{Disponibilidad NoSQL}
    NoSQL permite que el servidor de Madrid conteste ``OK'' de inmediato y ya se sincronizará con Tokio unos milisegundos después.
  \end{alertblock}
\end{frame}

\section{2. Sistemas Distribuidos}

\begin{frame}{2.1 ¿Qué es un Sistema Distribuido?}
    \textbf{Definición simple:} Es un grupo de ordenadores que trabajan juntos para que el usuario crea que está usando uno solo.
    
    \begin{itemize}
        \item \textbf{Nodo:} Cada servidor individual del sistema.
        \item \textbf{Cluster:} El conjunto de todos los nodos.
        \item \textbf{Red:} El ``pegamento'' que los une mediante mensajes (JSON, binario, etc.).
    \end{itemize}
    
    \begin{center}
    \begin{tikzpicture}[scale=0.6]
        \node (user) [cloud] {Cliente (App)};
        \node (n1) [node, below left=of user] {Nodo 1};
        \node (n2) [node, below=of user] {Nodo 2};
        \node (n3) [node, below right=of user] {Nodo 3};
        \draw [arrow] (user) -- (n1);
        \draw [arrow] (user) -- (n2);
        \draw [arrow] (user) -- (n3);
        \draw [dashed, myblue, <->] (n1) -- (n2);
        \draw [dashed, myblue, <->] (n2) -- (n3);
    \end{tikzpicture}
    \end{center}
    \note{En SQL solemos tener un servidor central (Monolito). En NoSQL, el cluster es la unidad básica.}
\end{frame}

\begin{frame}{2.2 Escalabilidad: Vertical vs Horizontal}
    \begin{columns}
        \begin{column}{0.5\textwidth}
            \textbf{Vertical (Scaling UP)}
            \begin{itemize}
                \item Comprar más RAM o CPU para el mismo servidor.
                \item \textbf{Problema:} Es carísimo y tiene un límite físico.
            \end{itemize}
        \end{column}
        \begin{column}{0.5\textwidth}
            \textbf{Horizontal (Scaling OUT)}
            \begin{itemize}
                \item Añadir más servidores baratos al grupo.
                \item \textbf{Ventaja:} Es la base de NoSQL y de la nube (Cloud).
            \end{itemize}
        \end{column}
    \end{columns}
    \begin{exampleblock}{La analogía del transporte}
        Vertical es comprar un camión más grande. Horizontal es contratar a 10 mensajeros en moto.
    \end{exampleblock}
\end{frame}

\begin{frame}{2.3 Conceptos de Red: Latencia y Partición}
    Para entender NoSQL hay que entender qué falla en la red:
    \begin{itemize}
        \item \textbf{Latencia:} El tiempo que tarda un mensaje en viajar. Nunca es cero.
        \item \textbf{Partición de Red:} Cuando dos nodos no pueden hablar entre sí (se rompe el cable, falla el switch).
    \end{itemize}
    \begin{alertblock}{La importancia del fallo}
        En un sistema de un solo servidor (SQL clásico), si el servidor falla, todo para. En un sistema distribuido, si un nodo falla, el resto debe seguir trabajando.
    \end{alertblock}
\end{frame}

\begin{frame}{2.4 Roles en el Cluster}
    No todos los nodos hacen lo mismo:
    \begin{itemize}
        \item \textbf{Primario / Maestro:} El que manda. Suele recibir las escrituras.
        \item \textbf{Secundario / Réplica:} Copia los datos del maestro. Sirve para lecturas rápidas o por si el maestro muere.
        \item \textbf{Árbitro:} Un nodo que no tiene datos, solo sirve para votar quién es el nuevo maestro si el actual falla.
    \end{itemize}
    \note{Esto es lo que verán los alumnos cuando configuren un 'Replica Set' en MongoDB.}
\end{frame}

\begin{frame}{2.5 Transparencia del Sistema}
    Un buen sistema distribuido debe ser transparente:
    \begin{itemize}
        \item \textbf{De ubicación:} Al programador le da igual en qué servidor esté el dato.
        \item \textbf{De fallo:} Si un nodo se desconecta, la aplicación ni se entera (el sistema redirige el tráfico).
        \item \textbf{De replicación:} El programador escribe una vez, el sistema lo copia a varios servidores.
    \end{itemize}
    \begin{exampleblock}{Ejemplo práctico}
        En MongoDB, tú haces un \texttt{db.users.insert()}. No sabes si el dato se ha guardado en el servidor de Madrid o de París, ni cuántas copias se han hecho. El SGBD lo gestiona por ti.
    \end{exampleblock}
\end{frame}

\section{3. El Agregado y el Modelado}

\begin{frame}{3.1 ¿Qué es un Agregado? (RA7-CE.c)}
  \textbf{Definición Profunda:} Es un grupo de objetos relacionados que queremos tratar como una \textbf{unidad indivisible}. 
  
  \begin{itemize}
    \item \textbf{SQL (La Despensa):} Tienes los huevos en un sitio, la harina en otro y la leche en la nevera. Para hacer una tarta (objeto), tienes que ir a buscarlos y mezclarlos (JOIN).
    \item \textbf{NoSQL (La Tarta ya hecha):} Guardas la tarta completa en una caja. Si quieres tarta, coges la caja y ya está. No pierdes tiempo buscando ingredientes.
  \end{itemize}
  
  \begin{exampleblock}{La unidad de Atomicidad}
    En NoSQL, la base de datos te garantiza que o se guarda el ``documento'' entero o no se guarda nada. No hay términos medios.
  \end{exampleblock}
  \note{Explicar que el agregado es el 'límite' de la transacción. En NoSQL no solemos tener transacciones entre documentos, por eso metemos todo lo importante dentro de uno solo.}
\end{frame}

\begin{frame}[fragile]{3.2 Ejemplo Real: El Carrito de la Compra}
  ¿Cómo vería un carrito una App de móvil?
  \begin{itemize}
    \item \textbf{En SQL:} Necesitarías la tabla \texttt{CARRITOS}, la tabla \texttt{LINEAS\_CARRITO} y la tabla \texttt{PRODUCTOS}. Tres consultas o un triple JOIN.
    \item \textbf{En NoSQL:} Es un único ``Agregado'' JSON.
  \end{itemize}
  \begin{lstlisting}[language=SQL]
{
  "cart_id": "ABC-123",
  "items": [
    { "id": "P01", "nombre": "Teclado", "precio": 85, "tags": ["gaming", "rgb"] },
    { "id": "P99", "nombre": "Cable", "precio": 12 }
  ],
  "total": 97,
  "ultima_modificacion": "2024-04-16T10:00:00Z"
}
  \end{lstlisting}
  \note{Fijaos en los 'tags'. En SQL eso sería OTRA tabla más. Aquí, el teclado ``lleva sus etiquetas puestas''.}
\end{frame}

\begin{frame}{3.3 Redundancia Controlada: El Snapshot}
  ¿Por qué repetimos el nombre y el precio en el carrito si ya están en el catálogo? 
  
  \begin{enumerate}
    \item \textbf{Velocidad (Performance):} Si el usuario abre el carrito 100 veces, le ahorramos a la base de datos 100 JOINs. Leer es casi instantáneo.
    \item \textbf{Consistencia Histórica (Snapshot):} Si mañana el teclado sube de 85€ a 100€, el carrito que el usuario llenó hoy \textbf{debe} mantener los 85€. Si solo tuviéramos un ID al producto, el precio cambiaría ``mágicamente'' en el carrito del cliente.
  \end{enumerate}
  
  \begin{alertblock}{La Redundancia es tu amiga}
    En NoSQL preferimos gastar un poco más de disco (que es barato) para comprar tiempo de CPU (que es caro) y coherencia en los datos.
  \end{alertblock}
\end{frame}

\section{4. CAP y PACELC}

\begin{frame}{4.1 Teorema CAP: La dura realidad}
  Imagina que tienes dos servidores (A y B) sincronizados. De repente, se corta el cable entre ellos (\textbf{Partición}). Un usuario quiere escribir un dato en A. ¿Qué haces?
  
  \begin{itemize}
    \item \textbf{Opción 1: Consistencia (C).} Bloqueas el servidor A y dices ``No puedo escribir hasta que vuelva la red''. El usuario se enfada porque el sistema no funciona, pero los datos son correctos.
    \item \textbf{Opción 2: Disponibilidad (A).} Dejas que el usuario escriba en A. El sistema funciona, pero el servidor B tiene datos viejos. Los datos no son iguales en todo el cluster.
  \end{itemize}
  
  \begin{exampleblock}{La elección obligatoria}
    Como la red (\textbf{P}) fallará algún día, tienes que elegir: ¿Prefieres que tu sistema sea \textbf{CP} (siempre correcto pero a veces caído) o \textbf{AP} (siempre funcionando pero a veces con datos viejos)?
  \end{exampleblock}
\end{frame}

\begin{frame}{4.2 PACELC: ¿Y si la red NO falla?}
    El teorema CAP solo nos preocupa cuando hay problemas. El profesor Daniel Abadi dijo: ``¿Y qué pasa el 99\% del tiempo cuando la red funciona bien?''.
    
    \begin{itemize}
        \item \textbf{PAC:} Si hay \textbf{P}artición, eliges entre \textbf{A}vailability o \textbf{C}onsistency.
        \item \textbf{ELC:} \textbf{E}lse (si no hay fallo), eliges entre \textbf{L}atency o \textbf{C}onsistency.
    \end{itemize}
    
    \begin{alertblock}{El precio de la perfección}
        Incluso con la red perfecta, si quieres que el Nodo B tenga el dato \textit{al mismo milisegundo} que el Nodo A, tienes que hacer esperar al usuario mientras el mensaje viaja. Esa espera es la \textbf{Latencia (L)}.
    \end{alertblock}
\end{frame}

\begin{frame}{4.3 El modelo BASE: El hermano relajado de ACID}
    En SQL usamos ACID (Consistencia absoluta). En NoSQL usamos \textbf{BASE}:
    
    \begin{itemize}
        \item \textbf{BA (Basically Available):} El sistema parece que siempre funciona.
        \item \textbf{S (Soft state):} El estado de los datos puede cambiar con el tiempo sin que el usuario haga nada (porque los nodos se están sincronizando de fondo).
        \item \textbf{E (Eventual consistency):} Al final, si dejas de escribir, todos los nodos tendrán lo mismo. Pero ``al final'' puede ser dentro de 200ms.
    \end{itemize}
    
    \note{Ejemplo para clase: Los 'likes' de una foto de Instagram. Si tú ves 100 y tu amigo ve 102, no pasa nada. Al cabo de un rato, ambos veréis 105. Es consistencia eventual.}
\end{frame}


% Sección 5 subdividida
\section{5a. Bases de Datos de Clave-Valor}

\begin{frame}{5.1 Concepto: El diccionario persistente}
    Es el modelo más sencillo y rápido que existe.
    \begin{itemize}
        \item \textbf{Estructura:} Un par \texttt{Llave} $\rightarrow$ \texttt{Valor}.
        \item \textbf{Opacidad:} Para la base de datos, el valor es una ``caja negra'' (BLOB). No sabe qué hay dentro, solo lo guarda y lo devuelve.
        \item \textbf{Analogía:} Es como las taquillas de un gimnasio. El número de taquilla es la \textbf{Llave} y lo que guardas dentro es el \textbf{Valor}. El gimnasio no sabe si guardas una mochila o unos zapatos.
    \end{itemize}
    \begin{exampleblock}{SGBD Referente: Redis}
        Es ultra-rápido porque guarda los datos en la memoria RAM, no en el disco duro.
    \end{exampleblock}
\end{frame}

\begin{frame}{5.2 Arquitectura: ¿Cómo escala?}
    ¿Cómo puede un cluster de 100 servidores saber dónde está una llave?
    \begin{itemize}
        \item \textbf{Hashing Consistente:} El sistema aplica una función matemática a la llave (Hash) para obtener un número.
        \item \textbf{El Anillo:} Ese número corresponde a un servidor específico en el cluster.
    \end{itemize}
    \begin{alertblock}{Velocidad O(1)}
        No importa si tienes 10 registros o 10 billones; el tiempo para encontrar una llave es siempre el mismo porque vas directo al servidor que la tiene.
    \end{alertblock}
    \note{Mencionar que aquí no hay JOINs ni escaneos de tabla. Es ``dame esto'' o ``guarda aquello''.}
\end{frame}

\begin{frame}{5.3 Casos de Uso: ¿Cuándo usarlas?}
    No sirven para todo, pero en lo suyo son imbatibles:
    \begin{itemize}
        \item \textbf{Gestión de Sesiones:} Guardar si el usuario ``Pepe'' está logueado en la web.
        \item \textbf{Carritos de compra volátiles:} Datos que cambian mucho y se borran pronto.
        \item \textbf{Caché de base de datos:} Si una consulta SQL tarda 2 segundos, guardas el resultado en Redis y la próxima vez tardará 1 milisegundo.
    \end{itemize}
    \begin{exampleblock}{Ejemplo real: Twitter}
        Usa Redis para guardar el 'Timeline' de los usuarios famosos y servirlo al instante a millones de personas.
    \end{exampleblock}
\end{frame}

\begin{frame}{5.4 Características: TTL y Tipos de Datos}
    Redis añade ``superpoderes'' al modelo básico:
    \begin{itemize}
        \item \textbf{TTL (Time To Live):} Puedes decirle: ``Guarda esta sesión, pero bórrala automáticamente en 30 minutos''. Es ideal para autolimpieza.
        \item \textbf{No solo Strings:} El valor puede ser una \textbf{Lista} (para colas de tareas), un \textbf{Set} (para etiquetas únicas) o un \textbf{Hash} (para mini-objetos).
    \end{itemize}
    \note{El TTL ahorra muchísimo código de mantenimiento. En SQL tendrías que hacer un 'DELETE FROM sesiones WHERE fecha < ...' cada hora.}
\end{frame}

\begin{frame}[fragile]{5.5 Ejemplo Práctico (Redis CLI)}
    Imagina que estamos en la terminal de Redis:
    \begin{lstlisting}[language=SQL]
# Guardar un token de usuario que expira en 1 hora (3600 seg)
SET session:user:45 "token_a78fb2" EX 3600

# Recuperar el token
GET session:user:45

# Usar una lista para una cola de mensajes (mensajeria)
LPUSH notificaciones:user:45 "Has recibido un mensaje"
LPUSH notificaciones:user:45 "Tu pedido ha salido"

# Sacar el ultimo mensaje (LIFO)
LPOP notificaciones:user:45
    \end{lstlisting}
\end{frame}

\section{5b. Bases de Datos Documentales}

\begin{frame}{5.1 Concepto: El documento autodescriptivo}
    A diferencia de SQL, aquí no guardamos ``filas'', guardamos \textbf{Documentos} (normalmente JSON o BSON).
    \begin{itemize}
        \item \textbf{Autodescriptivo:} Cada documento lleva sus propias ``columnas'' (llaves) dentro.
        \item \textbf{Jerárquico:} Puedes meter listas y otros objetos dentro de un documento.
        \item \textbf{SGBD Referente:} \textbf{MongoDB}.
    \end{itemize}
    \begin{exampleblock}{La analogía del Archivador}
        SQL es una hoja de Excel rígida. NoSQL Documental es una carpeta con folios: cada folio puede tener una estructura distinta, pero todos están guardados en el mismo cajón.
    \end{exampleblock}
\end{frame}

\begin{frame}{5.2 Esquema Flexible (Schema-less)}
    ¿Qué significa realmente que no hay esquema?
    \begin{itemize}
        \item \textbf{Polimorfismo:} En la colección \texttt{alumnos}, un alumno puede tener \texttt{telefono} y otro puede tener un array \texttt{telefonos: [``casa'', ``movil'']}. A la base de datos le da igual.
        \item \textbf{Evolución sin dolor:} Si tu App necesita un campo nuevo, lo añades en el código y listo. No hay que hacer un \texttt{ALTER TABLE} que bloquee la base de datos durante horas.
    \end{itemize}
    \begin{alertblock}{La responsabilidad es del programador}
        Como la base de datos no valida el esquema, es tu código el que debe asegurarse de que los datos tienen sentido.
    \end{alertblock}
\end{frame}

\begin{frame}{5.3 Modelado: ¿Embeber o Referenciar?}
    Esta es la pregunta del millón en NoSQL Documental:
    \begin{itemize}
        \item \textbf{Embeber (Embedding):} Metes los datos hijos dentro del padre.
        \begin{itemize}
            \item \textit{Cuándo:} Relaciones 1:1 o 1:pocos (ej: direcciones de un usuario).
            \item \textit{Ventaja:} Lees todo de un golpe (Sin JOINs).
        \end{itemize}
        \item \textbf{Referenciar (Linking):} Guardas el ID del hijo.
        \begin{itemize}
            \item \textit{Cuándo:} Relaciones 1:Muchos que crecen sin parar (ej: comentarios de un post viral).
            \item \textit{Ventaja:} Evitas que el documento supere el límite de tamaño (16MB en Mongo).
        \end{itemize}
    \end{itemize}
\end{frame}

\begin{frame}{5.4 Indexación y Potencia de Consulta}
    Aunque sea NoSQL, puedes buscar de forma muy avanzada:
    \begin{itemize}
        \item \textbf{Índices secundarios:} Puedes crear un índice sobre el email, la edad o incluso sobre un campo que está dentro de una lista.
        \item \textbf{Agregaciones:} Puedes filtrar, agrupar y calcular medias como en SQL, pero usando un ``pipeline'' (tubería) de etapas.
    \end{itemize}
    \begin{exampleblock}{Consultas Ad-hoc}
        A diferencia de Clave-Valor, aquí puedes decir: ``Dame todos los usuarios de Madrid que tengan más de 20 años y les guste el Rock''.
    \end{exampleblock}
\end{frame}

\begin{frame}[fragile]{5.5 Ejemplo Práctico (MongoDB Shell)}
    \begin{lstlisting}[language=SQL]
// Insertar un profesor con sus modulos embebidos
db.profesores.insertOne({
  nombre: "Vicente",
  especialidad: "Informatica",
  modulos: [
    { nombre: "Bases de Datos", horas: 6 },
    { nombre: "Sostenibilidad", horas: 3 }
  ],
  activo: true
})

// Buscar profesores que den mas de 5 horas de Bases de Datos
db.profesores.find({
  "modulos.nombre": "Bases de Datos",
  "modulos.horas": { $gt: 5 }
})
    \end{lstlisting}
\end{frame}

\section{5c. Bases de Datos de Familias de Columnas}

\begin{frame}{5.1 Concepto: El Mapa de Mapas}
    No pienses en tablas de Excel. Piensa en un \textbf{Mapa de Mapas} ordenado por una llave.
    \begin{itemize}
        \item \textbf{RowKey (Clave de Partición):} Es la llave que decide en qué servidor del cluster se guarda el dato físicamente.
        \item \textbf{Columna (Column):} Un par \texttt{Nombre: Valor} con marca de tiempo.
        \item \textbf{SGBD Referente:} \textbf{Apache Cassandra}.
    \end{itemize}
    \begin{exampleblock}{Diferencia con SQL}
        En SQL, una fila es una unidad física. En Cassandra, la fila es solo una agrupación lógica de columnas bajo una misma clave.
    \end{exampleblock}
\end{frame}

\begin{frame}{5.2 ¿Cómo guardan los datos? (LSM-Trees)}
    ¿Por qué Cassandra escribe más rápido que cualquier SQL?
    \begin{itemize}
        \item \textbf{Escritura Secuencial:} No ``busca'' el dato en el disco para editarlo (eso es lento). Simplemente escribe la nueva versión al final de un archivo de log.
        \item \textbf{Memtable:} Los datos se guardan en RAM y se vuelcan a disco en bloque.
    \end{itemize}
    \begin{alertblock}{No hay borrados reales}
        Si borras un dato, se escribe una marca llamada ``Tombstone'' (lápida). El dato se borrará físicamente más tarde durante una limpieza automática (Compaction).
    \end{alertblock}
\end{frame}

\begin{frame}{5.3 Tablas Dispersas (Sparse Tables)}
    En SQL, si una tabla tiene 100 columnas y solo rellenas 2, el resto son NULLs que ocupan sitio o lógica.
    \begin{itemize}
        \item \textbf{Sin NULLs:} En una BBDD de columnas, si una fila no tiene el dato ``Email'', la columna ``Email'' simplemente \textbf{no existe} para esa fila. No ocupa ni un bit.
        \item \textbf{Familias de Columnas:} Agrupamos columnas que suelen leerse juntas (ej: ``Datos Personales'' separada de ``Estadísticas de Uso'').
    \end{itemize}
    \note{Esto permite que el disco solo lea los trozos de archivo que contienen las columnas que has pedido, ignorando el resto.}
\end{frame}

\begin{frame}{5.4 Modelado: Query-First (Diseño por consulta)}
    ¡Aquí la normalización está prohibida!
    \begin{itemize}
        \item \textbf{Primero la pregunta, luego la tabla:} No diseñas tablas para guardar datos, diseñas tablas para responder preguntas.
        \item \textbf{Desnormalización total:} Si necesitas ver los usuarios por Ciudad y también por Edad, creas \textbf{dos tablas} con los mismos datos repetidos, optimizadas para cada búsqueda.
    \end{itemize}
    \begin{exampleblock}{La regla de oro}
        Un \texttt{SELECT} nunca debe requerir un \texttt{JOIN}. Si necesitas datos de dos sitios, guárdalos ya juntos.
    \end{exampleblock}
\end{frame}

\begin{frame}[fragile]{5.5 Ejemplo Visual de Almacenamiento}
    Imagina una tabla de usuarios en Cassandra:
    \begin{table}[]
    \tiny
    \begin{tabular}{lll}
    \toprule
    \textbf{RowKey (ID)} & \textbf{Family: Perfil} & \textbf{Family: Actividad} \\ \midrule
    user\_1 & \{nombre: ``Ana'', edad: 20\} & \{login: ``2024-04-16''\} \\
    user\_2 & \{nombre: ``Luis''\} & \{login: ``2024-04-15'', ip: ``10.0.0.1''\} \\
    \bottomrule
    \end{tabular}
    \end{table}
    \begin{itemize}
        \item \textbf{Dato Clave:} Luis no tiene campo 'edad'. Luis tiene un campo 'ip' que Ana no tiene.
        \item \textbf{Físico:} El disco guarda la familia ``Perfil'' en un archivo y ``Actividad'' en otro. Si solo quieres nombres, el disco no toca los datos de actividad.
    \end{itemize}
\end{frame}

\section{5d. Bases de Datos de Grafos}

\begin{frame}{5.1 Concepto: La relación es lo primero}
    En todas las BBDD anteriores, la relación era un ID guardado en algún sitio. Aquí la relación es un objeto con nombre y propiedades.
    \begin{itemize}
        \item \textbf{Nodos (Vértices):} Las entidades (ej: Persona, Ciudad, Producto).
        \item \textbf{Relaciones (Aristas):} Las conexiones (ej: CONOCE\_A, VIVE\_EN, COMPRÓ).
        \item \textbf{Propiedades:} Atributos de ambos (ej: Maria tiene 20 años, Maria conoce a Juan desde 2010).
        \item \textbf{SGBD Referente:} \textbf{Neo4j}.
    \end{itemize}
\end{frame}

\begin{frame}{5.2 Adyacencia sin índices (Index-free Adjacency)}
    ¿Por qué son tan rápidos buscando relaciones?
    \begin{itemize}
        \item \textbf{Punteros físicos:} Cada nodo tiene grabada la dirección de memoria de sus vecinos.
        \item \textbf{Sin búsquedas:} Para ir de ``Juan'' a sus ``Amigos'', el disco no busca en una tabla; simplemente salta a las direcciones de memoria que ya conoce.
    \end{itemize}
    \begin{exampleblock}{La analogía del laberinto}
        SQL es como un laberinto donde en cada esquina tienes que mirar un mapa (Índice). Grafos es como un laberinto donde hay un hilo de Ariadna que te lleva directo a la salida.
    \end{exampleblock}
\end{frame}

\begin{frame}{5.3 ¿Cuándo usarlas? (Conectividad)}
    Si tu pregunta empieza por ``¿Cómo están conectados...?'', usa Grafos.
    \begin{itemize}
        \item \textbf{Redes Sociales:} Amigos de amigos, gente que quizás conozcas.
        \item \textbf{Detección de Fraude:} ¿Hay un grupo de personas enviándose dinero en círculo para blanquearlo?
        \item \textbf{Recomendaciones:} ``Gente que compró este libro y vive en tu ciudad también compró este otro''.
        \item \textbf{Logística:} Calcular la ruta más rápida entre dos puntos.
    \end{itemize}
\end{frame}

\begin{frame}{5.4 Lenguaje de Consulta: Cypher}
    Es el lenguaje de Neo4j. Es muy intuitivo porque usa ``dibujos'' hechos con texto (ASCII Art).
    \begin{itemize}
        \item Los nodos van entre paréntesis: \texttt{(p:Persona)}
        \item Las relaciones van entre corchetes con flechas: \texttt{-[:CONOCE]->}
        \item Los filtros son declarativos, como en SQL.
    \end{itemize}
    \begin{exampleblock}{Ejemplo de patrón}
        \texttt{(juan)-[:AMIGO\_DE]->(maria)}
    \end{exampleblock}
\end{frame}

\begin{frame}[fragile]{5.5 Ejemplo Práctico (Cypher)}
    Imagina que queremos buscar recomendaciones de amigos:
    \begin{lstlisting}[language=SQL]
// Crear Maria y Juan y su relacion
CREATE (p1:Persona {nombre: "Maria"})
CREATE (p2:Persona {nombre: "Juan"})
CREATE (p1)-[:CONOCE {desde: 2020}]->(p2)

// Buscar "Amigos de mis amigos" que no sean mis amigos directos
MATCH (yo:Persona {nombre: "Juan"})-[:AMIGO_DE]->(amigo)-[:AMIGO_DE]->(fof)
WHERE NOT (yo)-[:AMIGO_DE]->(fof)
RETURN fof.nombre
    \end{lstlisting}
    \note{Fijaos lo fácil que es representar 'dos saltos' en la red con la flecha encadenada. En SQL esto sería un JOIN de la tabla amistades consigo misma.}
\end{frame}


\section{6. Conclusión y Futuro}

\begin{frame}{6.1 Persistencia Políglota: El mundo real}
    ¿Cuál es la mejor base de datos? \textbf{Ninguna y todas a la vez.}
    
    \begin{itemize}
        \item En una empresa moderna, se usa la base de datos adecuada para cada parte de la App.
        \item \textbf{SQL:} Para la pasarela de pagos y transacciones bancarias (ACID).
        \item \textbf{Redis:} Para las sesiones de usuario y caché de la web.
        \item \textbf{MongoDB:} Para el catálogo de productos y el perfil del usuario.
        \item \textbf{Neo4j:} Para el sistema de recomendaciones (``A otros usuarios también les gustó...'').
    \end{itemize}
    \begin{alertblock}{No seas fanático}
        Un buen arquitecto de software no se casa con una tecnología. Elige la herramienta que mejor resuelva el problema del cliente.
    \end{alertblock}
\end{frame}

\begin{frame}{6.2 Matriz de Selección (RA7-CE.b)}
    \begin{table}[]
    \centering
    \tiny
    \begin{tabular}{@{}llll@{}}
    \toprule
    \textbf{Tipo} & \textbf{¿Cuándo usarla?} & \textbf{Punto Fuerte} & \textbf{Escalado} \\ \midrule
    \textbf{Relacional} & Datos estructurados & Integridad (ACID) & Vertical \\
    \textbf{Clave-Valor} & Caché, Sesiones & Latencia mínima & Horizontal \\
    \textbf{Documental} & Apps Web, Catálogos & Flexibilidad & Horizontal \\
    \textbf{Columnares} & Big Data, Logs & Escritura masiva & Horizontal \\
    \textbf{Grafos} & Relaciones complejas & Conectividad & Complejo \\ \bottomrule
    \end{tabular}
    \end{table}
    \note{Resumir que NoSQL no vino a matar a SQL, sino a solucionar lo que SQL no podía hacer bien.}
\end{frame}

\begin{frame}{6.3 Resumen del RA7: ¿Qué hemos aprendido?}
    Al finalizar este tema, debes ser capaz de:
    \begin{enumerate}
        \item \textbf{Diferenciar} los modelos NoSQL del modelo relacional.
        \item \textbf{Identificar} qué tipo de base de datos conviene según el caso de uso.
        \item \textbf{Entender} el impacto de la red y la distribución (Teorema CAP).
        \item \textbf{Modelar} datos usando agregados y redundancia controlada.
    \end{enumerate}
    \begin{exampleblock}{Próximo paso: La Práctica}
        Instalaremos un cluster de MongoDB y aprenderemos a realizar consultas complejas y pipelines de agregación.
    \end{exampleblock}
\end{frame}

\begin{frame}{6.4 Reflexión Final}
    ``Las bases de datos relacionales te dicen cómo deben ser tus datos. Las bases de datos NoSQL te preguntan qué quieres hacer con ellos.''
    
    \vspace{0.8cm}
    \begin{center}
        \Large \textbf{¿Preguntas?}
    \end{center}
\end{frame}


\end{document}
